\documentclass{ctexart}
\usepackage{amsmath, amssymb}
\usepackage{geometry}
\geometry{a4paper, margin=2cm}

\title{\textbf{第6章 刚体力学} --- 知识点与公式}
\author{}
\date{}

\begin{document}
\maketitle

\section*{6.1 刚体运动的描述}

\subsection*{刚体和自由度}
\textbf{刚体：}形状和大小完全不变的物体，组内任意两质点间距离保持不变。

\textbf{自由度：}确定物体位置所需的独立坐标数。
\begin{itemize}
  \item 质点：3 个自由度
  \item 自由刚体：6 个自由度（平动 3 + 转动 3）
\end{itemize}

\subsection*{刚体的平动}
刚体内任一直线始终保持和自身平行。
平动特点：各质点运动情况相同，可归结为质心运动。

\subsection*{刚体绕定轴转动}
运动学方程：$\theta = f(t)$

角速度：$\omega = \dfrac{\mathrm{d}\theta}{\mathrm{d}t}$，单位 rad/s

角加速度：$\beta = \dfrac{\mathrm{d}\omega}{\mathrm{d}t} = \dfrac{\mathrm{d}^{2}\theta}{\mathrm{d}t^{2}}$，单位 rad/s$^{2}$

角量与线量的关系：
\[
s = r\theta,\quad v = r\omega,\quad a_\tau = r\beta,\quad a_n = r\omega^{2}
\]

矢量关系：
\[
\vec{v} = \vec{\omega} \times \vec{r},\quad
\vec{a} = \vec{\beta} \times \vec{r} + \vec{\omega} \times \vec{v}
\]

刚体上各点加速度大小：
\[
|\vec{a}| = r\sqrt{\beta^{2} + \omega^{4}},\quad
\tan\alpha = \frac{|\beta|}{\omega^{2}}
\]

\section*{6.2 力矩 刚体绕定轴转动的转动定律}

\subsection*{力矩}
力对点的力矩：$\vec{M}_O = \vec{r} \times \vec{F}$

力对定轴的力矩：$M_z = F_\tau r = F_\perp h$

力矩求和：对同一参考点用矢量和，对同一轴用代数和。

\textbf{注意：}合外力为零与合力矩为零不等价；刚体合内力矩为零。

\subsection*{转动惯量}
\[
J = \sum_i \Delta m_i r_i^{2}\quad\text{（离散）},\qquad
J = \int r^{2}\,\mathrm{d}m\quad\text{（连续）}
\]

单位：kg$\cdot$m$^{2}$

计算三要素：总质量、质量分布、转轴位置。

\textbf{常见转动惯量：}
\begin{itemize}
  \item 细杆绕端点：$J = \dfrac13 ML^{2}$
  \item 细杆绕中心：$J = \dfrac{1}{12} ML^{2}$
  \item 圆环绕中心轴：$J = MR^{2}$
  \item 圆盘绕中心轴：$J = \dfrac12 MR^{2}$
  \item 圆柱绕中心轴：$J = \dfrac12 MR^{2}$
  \item 空心圆柱：$J = \dfrac12 M(R_1^{2} + R_2^{2})$
\end{itemize}

\subsection*{平行轴定理}
\[
J_{z'} = J_z + ML^{2}
\]
其中 $J_z$ 为绕过质心轴的转动惯量，$L$ 为两平行轴间垂直距离。

\subsection*{垂直轴定理（薄板）}
\[
J_z = J_x + J_y
\]
其中 $x$、$y$ 轴在薄板平面内，$z$ 轴垂直于薄板。

\subsection*{转动定律}
\[
M_z = J\beta
\]
作用在刚体上的合外力矩等于转动惯量与角加速度的乘积。

对比：$M \to F$，$J \to m$，$\beta \to a$。

\section*{6.3 绕定轴转动刚体的动能}

\subsection*{转动动能}
\[
E_k = \frac12 J\omega^{2}
\]

\subsection*{力矩的功}
\[
\mathrm{d}A = M\,\mathrm{d}\theta,\qquad
A = \int_{\theta_1}^{\theta_2} M\,\mathrm{d}\theta
\]
功率：$P = M\omega$

\subsection*{转动动能定理}
\[
A = \frac12 J\omega_2^{2} - \frac12 J\omega_1^{2}
\]

\subsection*{刚体的机械能}
\[
E = \frac12 J\omega^{2} + mgh_c
\]
其中 $h_c$ 为质心高度。机械能守恒定律对包含刚体的系统仍成立。

\section*{6.4 刚体的动量矩和动量矩守恒定律}

\subsection*{刚体定轴转动的动量矩}
\[
L_z = J_z\omega
\]

\subsection*{动量矩定理}
微分形式：$M_z\,\mathrm{d}t = \mathrm{d}(J\omega)$
积分形式：
\[
\int_{t_1}^{t_2} M_z\,\mathrm{d}t = J\omega_2 - J\omega_1
\]

\subsection*{动量矩守恒定律}
当 $M_z = 0$ 时，$L = J\omega = \text{常矢量}$。

注：守恒条件为合外力矩为零，不是冲量矩为零。
动量矩守恒定律与动量守恒定律彼此独立。

\subsection*{进动}
高速自转的陀螺在重力矩作用下发生进动。

进动角速度：
\[
\Omega = \frac{|\vec{M}|}{|\vec{L}|\sin\theta} = \frac{|\vec{M}|}{J\omega\sin\theta} \propto \frac{1}{\omega}
\]
适用条件：高速自转（$\omega \gg \Omega$）。

\end{document}
